\documentclass[10pt,a4paper]{amsart}
\usepackage[margin=19mm]{geometry}
\usepackage{amsmath,amssymb,mathtools}
\usepackage[scr=rsfs,cal=euler]{mathalfa}
\usepackage{xfrac}
\usepackage[unicode,hidelinks,bookmarks=false]{hyperref}
\newcommand{\ie}{i.e.\ }
\newcommand{\cf}{cf.\ }
\newcommand{\resp}{respectively\ }
\newcommand{\Subsection}{\subsection}
\providecommand{\nobreakdash}{\nobreak\hbox{-}}
\setlength{\emergencystretch}{2em}
\multlinegap0pt
\usepackage{amssymb}
\usepackage[shortlabels]{enumitem}
\usepackage{xcolor}
\usepackage{mathtools}
\usepackage{tikz}
\usepackage[normalem]{ulem}
\usepackage{cancel}
\newtheorem{proposition}{Proposition}
\usepackage{cleveref}
\crefname{proposition}{Proposition}{Propositions}
\newcommand{\abs}[1]{\left\vert#1\right\vert}
\newcommand{\e}{\varepsilon}
\newcommand{\eps}{\e}
\renewcommand{\geq}{\geqslant}
\DeclareMathOperator{\id}{Id}
\renewcommand{\leq}{\leqslant}
\newcommand{\norm}[1]{\left\Vert#1\right\Vert}
\title{A group action approach to the Daugavet property}
\author{Sheldon Dantas, Helena del R\'{\i}o, Tomáš Raunig}
\date{Source: arXiv:2606.20429v1 (CC BY 4.0)}
\begin{document}
\maketitle

\noindent\emph{Source and licence.} A group action approach to the Daugavet property. Version arXiv:2606.20429v1 (CC BY 4.0), distributed by arXiv under the Creative Commons Attribution 4.0 International licence, http://creativecommons.org/licenses/by/4.0/, as recorded in the arXiv licence metadata for this exact version and shown on the abstract page https://arxiv.org/abs/2606.20429v1. Full licence notice: \texttt{licenses/arxiv-2606.20429-CC-BY-4.0.txt}.\par\smallskip

\noindent\textit{Editorial scope.} Counted unit: the article's proposition stability-ell1 (source0.tex lines 1223--1234). The article's complete original proof of it is delivered with it (source0.tex lines 1236--1349, one \texttt{proof} environment of 114 lines). No mathematics is written, repaired or re-derived: the statement and the proof are the article's own lines, in the article's own notation, and no retained line is substituted or re-flowed.  Retained uncounted as the source-local context the proof needs: lines 514--516.  Nothing was omitted from any retained span, so the package carries no omission record.  No retained line contains a citation, so the wrapper carries no bibliography and no bibliography entry.  No retained line contains a figure environment, an image-inclusion command or drawing code, so no image is delivered and the package declares no asset dependency.  Editorial formatting only, in addition: an AMS \texttt{amsart} preamble that restates the article's own declarations and macros used by the retained lines, namely the environments \texttt{proposition} on separate counters, together with the standard packages those lines need.  The retained environments print in this document as Proposition 1 in order of appearance, whereas the article prints the same items with the numbering of its own full document.  The complete native source of the article as distributed in the arXiv e-print for this exact version is bundled unchanged as \texttt{sources/203215/source0.tex}.
	\begin{proposition} \label{fact-DP-implies-GDP}
		Let $X$ be a $G$-Banach space and let $H \leq G$ be a~subgroup of $G$. If $X$ has the $H$-DPr, then $X$ has the $G$-DPr. 
	\end{proposition}

	\begin{proposition} \label{stability-ell1}
		Let $X_i$ be $G_i$-Banach spaces for $i \in I$, where $I$ is any index set, and define $X = \left(\sum_{i \in I} X_i\right)_{\ell_1}$. Consider the coordinatewise action given by
		\begin{equation*}
			(g_i)_{i \in I}(x_i)_{i \in I} = (g_i x_i)_{i \in I}
		\end{equation*}
		for every $(x_i)_{i \in I} \in X$ and $(g_i)_{i \in I} \in G$, where $G$ is defined below. Then, we have that the following are equivalent.
		\begin{itemize}
			\item[(a)] $X$ has the $G$-DPr for $G = \prod_{i\in I} G_i$.
			\item[(b)] $X$ has the $G$-DPr for $G = \bigoplus_{i \in I} G_i$, w.r.t. the restriction of the same action.
			\item[(c)] Each $X_i$ has the $G_i$-DPr.
		\end{itemize}
	\end{proposition}

	\begin{proof} (b) $\Rightarrow$ (a). Since the direct sum is a subgroup of the direct product, this implications follows from \Cref{fact-DP-implies-GDP}.  
		
		\noindent
		(a) $\Rightarrow$ (c). Assume that $X$ has the $G$-DPr and fix an arbitrary $k \in I$. We show that $X_k$ has the $G_k$-DPr. Let $T_k: X_k \to X_k$ be a rank-one operator and without loss of generality assume that $\norm{T_k} = 1$. Choose $\eps > 0$ arbitrarily. We define an extension $\hat{T}: X \to X$ by
		\[
		\hat{T}((x_i)_{i \in I}) = (0, \dots, 0, T_k(x_k), 0, \dots)
		\]
		where the non-zero term is in the $k$-th position. Then $\hat{T}$ is a rank-one operator on $X$ with $\|\hat{T}\| = \norm{T_k} = 1$. Since $X$ has the $G$-DPr, we have
		\[
		\sup_{g \in G} \norm{\id_X + g \circ \hat{T}} = 1 + \|\hat{T}\| = 2.
		\]
		So there are $g = (g_i)_{i \in I} \in G$ and $x = (x_i)_{i \in I} \in S_X$ such that
		\begin{equation*}
			\norm{x + g \circ \hat{T}(x)} > 2 - \eps.
		\end{equation*}
		Observe that
		\[
		(\id_X + g \circ \hat{T})(x) = x + g(\hat{T}x) = (x_1, \dots, x_k + g_k T_k x_k, \dots).
		\]
		Calculating the norm in $X = (\sum X_i)_{\ell_1}$, we get
		\begin{align*}
			2 - \eps
			&< \norm{(\id_X + g \circ \hat{T})(x)}
			= \sum_{i \neq k} \norm{x_i}_{X_i} + \norm{x_k + g_k T_k x_k}_{X_k}\\
			&= (1 - \norm{x_k}_{X_k}) + \norm{x_k + g_k T_k x_k}_{X_k}
			\leq (1 - \norm{x_k}_{X_k}) + \norm{x_k} (1 + \norm{T_k})\\
			&= 1 + \norm{x_k}_{X_k},
		\end{align*}
		implying that $\norm{x_k}_{X_k} > 1 - \eps$ and consequently $\sum_{i \neq k} \norm{x_i}_{X_i} < \eps$.
		Now, consider the vector $y \in X$ defined by the embedding of $x_k$ (i.e., $y_k = x_k$ and $y_i = 0$ for $i \neq k$). Note that $\norm{x - y}_X = \sum_{i \neq k} \norm{x_i} < \e$. By the triangle inequality, we get that
		\begin{align*}
			\norm{x_k + g_k T_k x_k}_{X_k}
			&= \norm{y + g \circ \hat{T}(y)}_X \\
			&\geq \norm{x + g \circ \hat{T}(x)}_X - \norm{(x - y) + g \circ \hat{T}(x - y)}_X.
		\end{align*}
		Since $\norm{\id + g \circ \hat{T}} \leq 2$, we have that
		\[
		\norm{(\id_{X_k} + g_k T_k)(x_k)}_{X_k} > (2 - \e) - 2\e = 2 - 3\e.
		\]
		Finally, let $u = x_k / \norm{x_k} \in S_{X_k}$. Since $\norm{x_k} \leq 1$, we get
		\[
		\norm{(\id_{X_k} + g_k T_k)(u)}_{X_k} = \frac{1}{\norm{x_k}} \norm{(\id_{X_k} + g_k T_k)(x_k)} > 2 - 3\e.
		\]
		Since $\e$ was arbitrary, $\sup_{h \in G_k} \norm{\id_{X_k} + h \circ T_k} = 2 = 1 + \norm{T_k}$, completing the proof.
		
		\noindent
		(c) $\Rightarrow$ (b). Assume that each $X_i$ has the $G_i$-DPr. Let $T: X \to X$ be a~rank-one operator and again without loss of generality assume that $\norm{T} = 1$. We can write $T = \phi \otimes y$ for some $\phi \in S_{X^*}$ and $y \in S_X$.
		Recall that $X^*$ is isometric to $(\sum X_i^*)_{\ell_\infty}$.
		Thus, we can write $\phi = (\phi_i)_{i \in I}$ with $\phi_i \in X_i^*$ and $\norm{\phi} = \sup_i \norm{\phi_i}_{X_i^*}$.
		Similarly, we can write $y = (y_i)_{i \in I}$ with $y_i \in X_i$ and $\norm{y} = \sum_i \norm{y_i}_{X_i}$.
		Let $\eps > 0$ and find $\eta \in (0, 1)$ such that $3\eta < \eps$. Choose an index $k \in I$ such that $\norm{\phi_k} \ge 1 - \eta = \norm{\phi} - \eta$.
		We need to handle two cases.
		
		\noindent
		\textbf{Case 1:} $\norm{y_k} < \eta$.
		
		In this case, we can find $u \in S_{X_k}$ such that $\phi_k(u) > 1 - \eta$. Then, define $g \in G$ to be the identity on each coordinate and $x \in S_X$ by embedding $u$ into the $k$-th coordinate (i.e., $x_k=u$ and $x_i=0$ for $i \neq k$). We compute
		\begin{align*}
			\norm{x + gTx}
			&= \norm{x + \phi_k(u)y} 
			= \sum_{j \neq k} \norm{\phi_k(u) y_j}_{X_j} + \norm{u + \phi_k(u) y_k}_{X_k}.
		\end{align*}
		For the first term, we use $\sum_{j \neq k} \norm{y_j} = 1 - \norm{y_k}$:
		\[
		\sum_{j \neq k} \norm{\phi_k(u) y_j} = \abs{\phi_k(u)} (1 - \norm{y_k}) \ge (1-\eta)(1-\eta) > 1 - 2\eta.
		\]
		For the second term, using the triangle inequality to get
		\[
		\norm{u + \phi_k(u) y_k} \ge \norm{u} - \abs{\phi_k(u)}\norm{y_k} > 1 - \eta.
		\]
		Summing these gives $\norm{x + gTx} > 2 - 3\eta > 2 - \eps$.
		
		
		\noindent
		\textbf{Case 2:} $\norm{y_k} \ge \eta$.
		
		Consider the operator $T_k = \phi_k \otimes y_k$ on $X_k$.
		Since $X_k$ has the $G_k$-DPr, there exist $g_k \in G_k$ and $u \in S_{X_k}$ such that
		\[
		\norm{u + g_k T_k u}_{X_k} \ge 1 + \norm{T_k} - \eta^2.
		\]
		Substituting $T_k u = \phi_k(u)y_k$, this becomes
		\begin{equation} \label{eq:comp_k}
			\norm{u + g_k \phi_k(u) y_k}_{X_k}
			\ge 1 + \norm{\phi_k}\norm{y_k} - \eta^2
			\ge1 + \norm{\phi_k}\norm{y_k} - \eta.
		\end{equation}
		Using the triangle inequality, we obtain $\norm{g_k\phi_k(u)y_k} \ge 1 + \norm{\phi_k}\norm{y_k} - \eta^2 - 1$, which, using the fact that $G$ acts by isometries and dividing by $\norm{y_k}$, implies
		\begin{equation*}
			\abs{\phi_k(u)}
			\ge \norm{\phi_k} - \frac{\eta^2}{\norm{y_k}}
			\ge \norm{\phi_k} - \eta
			\ge 1 - 2\eta.
		\end{equation*}
		
		Now, define $x \in S_X$ by embedding $u$ into the $k$-th coordinate (i.e., $x_k=u$ and $x_i=0$ for $i \neq k$). Define $g \in G$ such that its $k$-th component is $g_k$ and $g_j$ is the identity element of $G_j$ for $j \neq k$. We compute
		\begin{align*}
			\norm{x + gTx}
			&= \norm{x + g(\phi(x)y)} 
			= \norm{x + g(\phi_k(u)y)} \\
			&= \sum_{j \neq k} \norm{\phi_k(u) y_j}_{X_j} + \norm{u + g_k \phi_k(u) y_k}_{X_k}.
		\end{align*}
		Thus, using \eqref{eq:comp_k} above, we get that
		\begin{align*}
			\norm{x + gTx} 
			&\ge |\phi_k(u)| \sum_{j \neq k} \norm{y_j} + (1 + \norm{\phi_k}\norm{y_k} - \eta)\\
			&\ge (1-2\eta) \sum_{j \neq k} \norm{y_j} + 1 + (1-\eta)\norm{y_k} - \eta\\
			&= 1 + \sum_{j \neq k} \norm{y_j} + \norm{y_k}
			- \left(2\sum_{j \neq k} \norm{y_j} + \norm{y_k} + 1\right)\eta\\
			&\ge 1 + \norm{y} - (2\norm{y} + 1) \eta\\
			&= 2 - 3\eta \ge 2-\eps.
		\end{align*}
		As $\e$ was chosen arbitrarily, $\sup_{g \in G} \norm{\id + g \circ T} \ge 1 + \norm{T}$ and $X$ has the $G$-DPr.
	\end{proof}

\end{document}
